% Part: lambda-calculus
% Chapter: introduction
% Section: syntax

\documentclass[../../../include/open-logic-section]{subfiles}

\begin{document}

\olfileid[pt]{lam}{int}{syn}
\olsection{A sintaxe do cálculo lambda}

Começamos com uma sequência de variáveis $x$, $y$, $z$,~\dots e alguns
símbolos de constante $a$, $b$, $c$,~\dots. O conjunto dos termos é definido
indutivamente da seguinte maneira:
\begin{enumerate}
\item Cada variável é um termo.
\item Cada constante é um termo.
\item Se $M$ e $N$ são termos, então $(MN)$ também é um termo.
\item Se $M$ é um termo e $x$ é uma variável, então $(\lambd[x][M])$ é um
  termo.
\end{enumerate}
Termos da forma $(MN)$ são chamados de \emph{aplicações}, e os da forma
$(\lambd[x][M])$, de \emph{abstrações}.

O sistema sem constante alguma é chamado de cálculo lambda \emph{puro}.
Trabalharemos principalmente no cálculo-$\lambd$ puro, de modo que todas as
letras minúsculas representarão variáveis. Usamos letras maiúsculas ($M$,
$N$ etc.) para representar termos do cálculo-$\lambd$.

Seguiremos algumas convenções de notação:

\begin{conv}
\begin{enumerate}
\item Quando os parênteses são omitidos, a aplicação é feita da esquerda
  para a direita. Por exemplo, se $M$, $N$, $P$ e $Q$ são termos, então
  $MNPQ$ abrevia $(((MN)P)Q)$.
\item Novamente, quando os parênteses são omitidos, a abstração lambda recebe
  o maior escopo possível. Por exemplo, $\lambd[x][MNP]$ é lido como
  $(\lambd[x][((MN)P)])$.
\item Um lambda pode ser usado para abstrair várias variáveis. Por exemplo,
  $\lambd[xyz][M]$ abrevia
  $\lambd[x][\lambd[y][\lambd[z][M]]]$.
\end{enumerate}
\end{conv}

Por exemplo,
\[
\lambd[xy][xxyx \lambd[z][xz]]
\]
abrevia
\[
\lambd[x][\lambd[y][((((xx)y)x)(\lambd[z][(xz)]))]].
\]
Você deve memorizar essas convenções. No início elas o deixarão maluco, mas
você se acostumará e, depois de algum tempo, elas o deixarão menos maluco do
que ter de lidar com um emaranhado de parênteses.

Dois termos que diferem apenas nos nomes das variáveis ligadas são chamados
de $\alpha$-equivalentes; por exemplo, $\lambd[x][x]$ e $\lambd[y][y]$.
Será conveniente considerá-los o ``mesmo'' termo; em outras palavras,
quando dizemos que $M$ e $N$ são iguais, também queremos dizer ``a menos de
renomeações das variáveis ligadas''. Variáveis que estão no escopo de um
$\lambd$ são chamadas de ``ligadas'', enquanto as demais são chamadas de
``livres''. Não há variáveis livres no exemplo anterior; mas, em
\[
(\lambd[z][yz])x
\]
$y$ e $x$ são livres, e $z$ é ligada.

\end{document}

